\documentclass[11pt,a4paper]{article}
\usepackage[margin=22mm]{geometry}
\usepackage{amsmath,amssymb,bm}
\usepackage{booktabs,longtable,array}
\usepackage{graphicx}
\usepackage{xcolor}
\usepackage{microtype}
\usepackage{hyperref}
\usepackage{listings}
\usepackage{caption}
\usepackage{subcaption}
\usepackage{fancyhdr}

\hypersetup{colorlinks=true,linkcolor=blue!55!black,urlcolor=blue!55!black,citecolor=blue!55!black}
\graphicspath{{./}}
\setlength{\parindent}{0pt}
\setlength{\parskip}{0.55em}
\pagestyle{fancy}
\setlength{\headheight}{14pt}
\fancyhf{}
\lhead{Cartesian-to-RASCAS static-sphere validation}
\rhead{\thepage}
\lstset{basicstyle=\ttfamily\small,breaklines=true,frame=single,columns=fullflexible}

\newcommand{\lya}{Ly$\alpha$}
\newcommand{\tauzero}{\tau_0}
\newcommand{\av}{a_{\rm v}}
\newcommand{\nhi}{n_{\rm HI}}
\newcommand{\kms}{\mathrm{km\,s^{-1}}}

\title{Detailed Validation of a Cartesian-to-RAMSES-Octree-to-RASCAS Pipeline\\
against the Analytical Homogeneous Static-Sphere \lya\ Solution}
\author{Reproducible workflow note}
\date{13 August 2026}

\begin{document}
\maketitle

\begin{abstract}
I construct a homogeneous, static, neutral-hydrogen sphere at $T=100$ K in a
$64^3$ Cartesian array, convert it to the RAMSES octree binary format, create a
spherical RASCAS domain of radius $0.3$ box units, emit monochromatic central
\lya\ photon packets, and compare the escaped spectra with the analytical
static-sphere solution.  I test centre-to-surface optical depths
$\tauzero=10^5,10^6,$ and $10^7$.  The pipeline is validated independently at
the Cartesian, RAMSES, RASCAS-domain, photon-input, and escaped-photon stages.
The peak locations, double-peaked morphology, normalization, and red--blue
symmetry agree with theory.  The normalized full-profile $L_1$ difference
improves from $0.1124$ at $\tauzero=10^5$ (one million photons) to $0.0649$ at
$10^6$ and $0.0446$ at $10^7$ (one hundred thousand photons each), consistent
with the increasing validity of the analytical diffusion approximation as
$\av\tauzero$ increases.  This document explains every physical choice,
conversion factor, binary interface, calculation, plot, discrepancy, and
acceptance decision step by step.
\end{abstract}

\tableofcontents
\newpage

\section{Purpose and validation philosophy}

My aim is not merely to produce a double-peaked spectrum.  I want to demonstrate
that the same physical problem survives the complete chain
\begin{align*}
\text{physical target}&\longrightarrow \text{Cartesian NumPy fields}
\longrightarrow \text{RAMSES octree files}\\
&\longrightarrow \text{RASCAS spherical domain}
\longrightarrow \text{Monte Carlo photons}
\longrightarrow \text{escaped spectrum}.
\end{align*}
At each arrow I perform an independent check.  This is important because a
plausible spectrum could still be produced with a wrong unit conversion, a
reversed ionization field, an incorrect field index, a temperature error, or a
diameter optical depth used where the analytical formula expects a radius
optical depth.

I use three optical-depth cases.  The $10^5$ case is a useful diagnostic but is
not deep in the asymptotic diffusion regime.  The $10^6$ and $10^7$ cases are
the stronger analytical tests.  For the initial controlled comparison I use
$10^5$ photon packets in each optical-depth case.  I additionally run
$10^6$ packets at $\tauzero=10^5$ to show which differences are Monte Carlo
noise and which are systematic profile-shape differences.

\section{Definition of the physical experiment}

\subsection{Geometry and gas state}

I choose a homogeneous sphere centred at
\[
\bm{x}_{\rm c}=(0.5,0.5,0.5)
\]
inside a unit computational box.  The spherical boundary radius is
\[
R=0.3\,L_{\rm box}.
\]
The code length unit is one kiloparsec,
\[
L_{\rm box}=\mathrm{unit\_l}=3.085677581282\times10^{21}\ \mathrm{cm},
\]
and therefore
\[
R=9.257032743846\times10^{20}\ \mathrm{cm}=0.3\ \mathrm{kpc}.
\]

I set $T=100$ K and use primordial hydrogen mass fraction $X_{\rm H}=0.76$.
The gas has zero bulk velocity, zero turbulent velocity, zero metallicity,
zero dust, and zero ionized fraction.  Cosmological expansion, recoil, and
helium-ion fields are absent.  The source is a point at the sphere centre,
with zero source velocity, emitting monochromatic photons at
$\lambda_\alpha=1215.67$ \AA.  Scattering is isotropic.

These choices isolate the basic resonant-transfer physics.  With no preferred
direction or velocity, the escaped spectrum must be symmetric about line
centre.  Any persistent red--blue asymmetry would be evidence of a numerical,
geometrical, or interface problem.

\subsection{Optical-depth convention}

I define the line-centre optical depth from the centre to the spherical surface:
\begin{equation}
\tauzero=\nhi\,\sigma_0(T)\,R.
\label{eq:tau}
\end{equation}
It is not the optical depth across the diameter.  This convention must match
the analytical solution and the RASCAS escape boundary.  Using $2R$ by mistake
would shift the expected peaks by $2^{1/3}$ and would be a genuine physics
failure.

\section{Atomic constants and thermodynamic quantities}

\subsection{RASCAS-compatible line quantities}

I calculate the constants from the installed RASCAS atomic definition rather
than inserting rounded textbook values.  The adopted values are
\[
m_{\rm ion}=1.008\,m_{\rm u},\qquad f_{12}=0.416,\qquad
A_{21}=6.265\times10^8\ \mathrm{s^{-1}}.
\]
At temperature $T$, the hydrogen thermal speed and Doppler width are
\begin{align}
v_{\rm th} &= \sqrt{\frac{2k_{\rm B}T}{m_{\rm ion}}},\\
\Delta\nu_D &= \frac{v_{\rm th}}{\lambda_\alpha}.
\end{align}
At 100 K I obtain
\[
v_{\rm th}=1.2844045175\ \kms,
\qquad
\Delta\nu_D=1.0565404406\times10^{10}\ \mathrm{Hz}.
\]
The Voigt damping parameter is
\begin{equation}
\av=\frac{A_{21}}{4\pi\Delta\nu_D}
=4.7187295447\times10^{-3}.
\end{equation}

The line-centre cross-section is evaluated with the same COLT approximation
to $H(\av,0)$ that is used by the installed RASCAS source:
\begin{equation}
\sigma_0 = \frac{\sqrt{\pi}e^2 f_{12}}{m_e c\,\Delta\nu_D}
H(\av,0)
=5.8642983367\times10^{-13}\ \mathrm{cm^2}.
\end{equation}
Matching this convention matters because a small cross-section mismatch becomes
a direct optical-depth mismatch.

\subsection{Temperature in code units}

The velocity and density units are
\[
\mathrm{unit\_v}=10^5\ \mathrm{cm\,s^{-1}},\qquad
\mathrm{unit\_d}=6.7699139342702\times10^{-23}\ \mathrm{g\,cm^{-3}}.
\]
For neutral primordial gas I use
\begin{equation}
\mu=\frac{1}{X_{\rm H}+(1-X_{\rm H})/4}=1.2195121951.
\end{equation}
The RAMSES temperature conversion is
\begin{equation}
T=\frac{P}{\rho}\,\mu\frac{m_p}{k_B}\,\mathrm{unit\_v}^2.
\end{equation}
Therefore the required pressure-to-density ratio at 100 K is
\[
\frac{P}{\rho}=0.6768607801.
\]
For $\gamma=5/3$ and $P=(\gamma-1)\rho u$, the internal energy field is
\begin{equation}
u=\frac{P/\rho}{\gamma-1}=1.0152911701.
\label{eq:u}
\end{equation}

\section{Choosing the density for each optical depth}

I rearrange \ref{eq:tau} to obtain
\begin{equation}
\nhi=\frac{\tauzero}{\sigma_0 R}.
\end{equation}
The corresponding RAMSES dimensionless mass density is
\begin{equation}
\rho_{\rm RAMSES}=\frac{\nhi m_p}{X_{\rm H}\,\mathrm{unit\_d}}.
\label{eq:rhoramses}
\end{equation}

\subsection{Mandatory converter density multiplier}

The Cartesian converter does not write its input density unchanged.  It applies
the built-in multiplier
\begin{equation}
F_\rho=\frac{h_0^2}{a_{\rm exp}^3}
=\frac{0.6736^2}{0.125^3}
=232.31332352.
\label{eq:factor}
\end{equation}
Although this factor has a cosmological form, I leave the converter unchanged.
I compensate for it in the input array:
\begin{equation}
\rho_{\rm Cartesian}=\frac{\rho_{\rm RAMSES}}{F_\rho}.
\label{eq:cartdensity}
\end{equation}
If I inserted the desired RAMSES density directly into the NumPy array, the
converted gas would be too dense and $\tauzero$ would be too large by a factor
of approximately 232.3.

\begin{table}[ht]
\centering
\caption{Densities derived without fitting or post-conversion rescaling.}
\label{tab:density}
\begin{tabular}{rrrr}
\toprule
$\tauzero$ & $\nhi$ [$\mathrm{cm^{-3}}$] & $\rho_{\rm Cartesian}$ & $\rho_{\rm RAMSES}$\\
\midrule
$10^5$ & $1.8420955095\times10^{-4}$ & $2.5777397642\times10^{-8}$ & $5.9884329180\times10^{-6}$\\
$10^6$ & $1.8420955095\times10^{-3}$ & $2.5777397642\times10^{-7}$ & $5.9884329180\times10^{-5}$\\
$10^7$ & $1.8420955095\times10^{-2}$ & $2.5777397642\times10^{-6}$ & $5.9884329180\times10^{-4}$\\
\bottomrule
\end{tabular}
\end{table}

\section{Step 1: constructing the Cartesian fields}

For every optical depth I make C-contiguous NumPy \texttt{float64} arrays of
shape $(64,64,64)$.  I use the axis order \texttt{array[x,y,z]}, which is the
order sampled by the converter.  I fill the complete cube homogeneously; I do
not apply a spherical mask at this stage.  The sphere is imposed later as the
exact RASCAS computational boundary.

The fields are
\begin{center}
\begin{tabular}{ll}
\toprule
File & Assigned value\\
\midrule
\texttt{density.npy} & case-dependent value in \ref{tab:density}\\
\texttt{velocity\_x.npy}, \texttt{velocity\_y.npy}, \texttt{velocity\_z.npy} & exactly zero\\
\texttt{metallicity.npy} & exactly zero\\
\texttt{internal\_energy.npy} & $1.0152911701$\\
\texttt{scalar\_01.npy} & $x_{\rm HII}=0$ exactly\\
\texttt{pressure.npy} & diagnostic only; ignored by the converter\\
\bottomrule
\end{tabular}
\end{center}

The converter derives pressure internally as
\begin{equation}
P=(\gamma-1)\rho u,
\end{equation}
so the saved pressure array is only a human-readable diagnostic.  After writing
each field I reload it from disk and check shape, dtype, C-contiguity, finiteness,
homogeneity, exact zeros, and min/max/mean.  Central-slice images are constant,
as required for a homogeneous experiment.

\section{Step 2: converting the Cartesian cube to a RAMSES octree}

I pass a validation-owned JSON configuration with \texttt{--config}; I do not
modify the converter's original configuration.  I run only the AMR and hydro
writers because particles are neither present nor required by
\texttt{CreateDomDump}.

The AMR metadata uses
\[
\texttt{levelmin}=\texttt{levelmax}=6,\qquad \texttt{boxlen}=1.
\]
The converter records the normal hierarchy through levels 1--6, but every
physical leaf is at level 6.  A level-6 uniform mesh has
\[
2^6\times2^6\times2^6=64^3=262,144
\]
leaf cells.  This exact count is recovered by the trusted RAMSES reader and
\texttt{yt} for all three cases.

The hydro field order written and later read is density, three velocity
components, pressure, metallicity, internal energy, and
\texttt{scalar\_01}.  I explicitly verify the field order because using
\texttt{itemp}, \texttt{imetal}, or \texttt{ihii} with the wrong positions
could produce a physically plausible but incorrect RASCAS mesh.

The converted validations recover
\begin{itemize}
\item exactly 262,144 level-6 leaves;
\item homogeneous case-dependent density after the factor in \ref{eq:factor};
\item $u=1.0152911701$ and the derived pressure;
\item $T=100$ K to floating-point accuracy;
\item exactly zero velocity, metallicity, and \texttt{scalar\_01}.
\end{itemize}

\section{Step 3: creating and validating the RASCAS spherical domain}

I run \texttt{CreateDomDump} with a sphere centred at $(0.5,0.5,0.5)$ and
radius 0.3.  The important reader settings are
\[
\texttt{ramses\_rt=T},\quad \texttt{itemp=5},\quad
\texttt{imetal=6},\quad \texttt{ihii=8},
\]
with absent helium ion fields set to zero, \texttt{cosmo=F},
\texttt{vturb\_kms=0}, and \texttt{ignoreDust=T}.

The mesh dump does not retain $x_{\rm HII}$ itself.  I therefore validate
$\texttt{scalar\_01}=0$ upstream in RAMSES, then check downstream that the
dumped $\nhi$ equals the expected total hydrogen density.  Each domain contains
33,696 selected level-6 leaves in 5,553 octs.

I calculate columns along the six Cartesian axes and 32 reproducibly seeded
random rays.  Every ray begins at the source and is clipped at the exact
spherical boundary.  For a uniform sphere,
\begin{equation}
N_{\rm HI}=\int_0^R \nhi\,ds=\nhi R,
\qquad
\tauzero^{\rm measured}=\sigma_0 N_{\rm HI}.
\end{equation}
All 38 rays recover the requested optical depth to floating-point precision,
without density fitting or rescaling.  The dumped temperature is 100 K and
speed, turbulence, and dust density are exactly zero.

\section{Step 4: constructing and validating photon packets}

I use \texttt{PhotonsFromSourceModel} to create unformatted sequential Fortran
files.  Every packet has
\begin{itemize}
\item source position $(0.5,0.5,0.5)$;
\item source velocity $(0,0,0)$;
\item an isotropically sampled unit direction;
\item wavelength 1215.67 \AA, corresponding to a single finite line-centre
frequency;
\item deterministic global seed $-20260813$ and an individual packet seed.
\end{itemize}

The photon input file contains records in this order: packet count, represented
real-photon rate, global seed, packet IDs, frequencies, positions, directions,
packet seeds, and source velocities.  My streaming validator checks the leading
and trailing Fortran record markers and supports logical records larger than
2 GiB.  It confirms the exact packet count, sequential unique IDs, central
positions, zero velocities, monochromatic finite frequencies, and unit direction
norms.

\section{Step 5: RASCAS transport assumptions}

RASCAS reads the validated spherical domain and photon file.  I use isotropic
HI scattering, no recoil, no dust, no bulk velocity, and no turbulent velocity.
I set \texttt{nDirections=0}; therefore the reported spectrum is made directly
from all successfully escaped packets and has no viewing-cone selection.
Because the problem is spherical and static, this is the natural angle-integrated
comparison with the analytical solution.

Core skipping is enabled with \texttt{xcritmax=3}.  This accelerates the very
large number of nearly local line-core scatterings and makes the high-optical-
depth experiments practical.  It is an explicitly documented approximation.
It should preserve the main escape spectrum when used appropriately, but at the
lowest $\av\tauzero$ it can be one contributor to small trough, width, or wing
differences.  I do not treat core skipping as a proven explanation for every
residual; the main known limitation at $\tauzero=10^5$ is the asymptotic nature
of the analytical expression itself.

\section{Reading the RASCAS result binary}

The escaped result \texttt{Cont.res} is also an unformatted sequential Fortran
file.  I read the following records in their written order:
\begin{enumerate}
\item total packet count;
\item packet IDs;
\item terminal status flags;
\item final positions;
\item external-frame frequencies $\nu$;
\item final propagation directions;
\item scattering counts;
\item propagation times.
\end{enumerate}
Only \texttt{status=1} packets enter the escaped spectrum.  I preserve and
report all status counts, so lost or unfinished packets cannot disappear
silently.  Every packet in every completed comparison reported here has status 1.

\section{Analytical static-sphere spectrum}

\subsection{Dimensionless frequency}

I convert each escaped external-frame frequency to Doppler units:
\begin{equation}
x=\frac{\nu-\nu_\alpha}{\Delta\nu_D},
\qquad
\nu_\alpha=\frac{c}{\lambda_\alpha}.
\label{eq:x}
\end{equation}
With this frequency convention, $x<0$ is red and $x>0$ is blue.  A wavelength
offset would have the opposite sign near line centre, so I never mix the two
conventions.

\subsection{Analytical equation and normalization}

For a homogeneous static sphere I evaluate
\begin{equation}
J(x)=\sqrt{\frac{\pi}{6}}\frac{1}{4\av\tauzero}
\frac{x^2}{1+\cosh\!\left[
\sqrt{\frac{2\pi^3}{27}}\frac{|x|^3}{\av\tauzero}
\right]}.
\label{eq:analytic}
\end{equation}
The convention in \ref{eq:analytic} satisfies
\[
\int_{-\infty}^{\infty}J(x)\,dx=\frac{1}{4\pi}.
\]
My Monte Carlo histogram is a unit-area probability density.  The directly
comparable analytical curve is consequently
\begin{equation}
p_{\rm analytic}(x)=4\pi J(x),\qquad
\int p_{\rm analytic}(x)\,dx=1.
\end{equation}
I evaluate the full curve at the histogram-bin centres.  I normalize its sampled
finite grid only for negligible quadrature truncation; I never change its width,
optical depth, or amplitude to fit RASCAS.

The approximate peak position is
\begin{equation}
|x_{\rm peak}|\simeq0.931(\av\tauzero)^{1/3}.
\label{eq:peak}
\end{equation}
Thus increasing optical depth forces photons farther into the wings before
escape and increases the peak separation approximately as $\tauzero^{1/3}$.

\section{Constructing the numerical spectrum and uncertainties}

I histogram the escaped $x$ values over $-100\le x\le100$ with bin width
$\Delta x=0.5$.  If bin $i$ contains $N_i$ packets and the total escaped count
is $N_{\rm esc}$, I calculate
\begin{equation}
p_{{\rm MC},i}=\frac{N_i}{N_{\rm esc}\Delta x},
\qquad
\sigma_i=\frac{\sqrt{N_i}}{N_{\rm esc}\Delta x}.
\label{eq:hist}
\end{equation}
Therefore
\[
\sum_i p_{{\rm MC},i}\Delta x=1.
\]
The raw count is always retained.  A smooth normalized curve alone can hide
whether a tail bin contains ten photons or ten thousand photons.

My main full-profile statistic is
\begin{equation}
L_1=\sum_i|p_{{\rm MC},i}-p_{{\rm analytic},i}|\Delta x.
\label{eq:l1}
\end{equation}
I also measure peak positions, peak separation, integrated red/blue packet
ratio, symmetry fraction
\[
S=\frac{|N_{\rm red}-N_{\rm blue}|}{N_{\rm red}+N_{\rm blue}},
\]
central probability $P(|x|<1)$, RMS density difference, mean scattering count,
and a chi-square-like diagnostic.  The latter is not interpreted as a formal
$p$-value because the normalized bins are correlated and the analytical
diffusion solution is approximate.

\section{Meaning of every generated plot}

\begin{description}
\item[Linear spectrum overlay]
This is the most direct physical comparison.  It emphasizes peak positions,
peak heights, widths, the central trough, and the probability-carrying portion
of the spectrum.  Agreement here is more important than visual agreement at
probabilities many orders of magnitude below the peak.

\item[Logarithmic spectrum overlay]
This displays the weak wings and rare events.  It is valuable for finding an
unexpected long tail, but it visually magnifies tiny absolute differences.
Bins with zero packets are plotted at a small positive floor so they can appear
on logarithmic axes; vertical drops to that floor are plotting artifacts, not
negative probability.  Isolated high tail bins can contain only one or a few
photons.

\item[Folded red and blue halves]
I reflect the red half about $x=0$ and compare it with the blue half versus
$|x|$.  This removes the common double-peaked shape and directly exposes
directional or frequency-sign asymmetry.  A static sphere should agree within
Monte Carlo noise.

\item[Absolute residual]
This shows $p_{\rm MC}-p_{\rm analytic}$.  Positive residuals identify where
RASCAS assigns more probability, and negative residuals where it assigns less.
It reveals whether a lower peak is compensated by broader shoulders or wings.

\item[Fractional residual]
This shows $(p_{\rm MC}-p_{\rm analytic})/p_{\rm analytic}$ only where the
analytical curve exceeds $10^{-4}$ of its maximum.  Without this mask,
division by a vanishing analytical tail would create huge but meaningless
ratios.

\item[Poisson count plot]
This displays $N_i$ with $\sqrt{N_i}$ error bars.  It is the clearest view of
the actual Monte Carlo information content and distinguishes statistically
meaningful discrepancies from sparsely populated tail fluctuations.
\end{description}

\section{Physical origin of the double-peaked spectrum}

A photon emitted at line centre sees the largest HI cross-section and therefore
has a very short mean free path.  Repeated resonant scatterings produce coupled
random walks in position and frequency.  Thermal atomic velocities move a
photon toward either frequency wing.  In a wing the cross-section is smaller,
the mean free path is longer, and a sufficiently large frequency excursion can
carry the photon to the boundary.  Because the gas is static and symmetric,
red and blue excursions are equally likely.

The expected signatures are consequently a suppressed line centre, two peaks,
equal integrated red and blue flux, and increasing peak separation with
$(\av\tauzero)^{1/3}$.  The simulation reproduces every one of these signatures.

\section{Quantitative results}

\begin{table}[ht]
\centering
\caption{Angle-integrated RASCAS comparisons.  The first three rows use
$10^5$ packets; the fourth is the $\tauzero=10^5$ photon-count refinement.}
\label{tab:results}
\resizebox{\textwidth}{!}{%
\begin{tabular}{rrrrrrrrrr}
\toprule
$\tauzero$ & $N_\gamma$ & $\av\tauzero$ & analytic $|x_p|$ & RASCAS $x_{p,r}$ & RASCAS $x_{p,b}$ & separation & red/blue & $S$ & $L_1$\\
\midrule
$10^5$ & $10^5$ & 471.873 & 7.2481 & $-7.25$ & 6.75 & 14.0 & 1.00977 & 0.00486 & 0.11329\\
$10^6$ & $10^5$ & 4718.730 & 15.6156 & $-15.25$ & 15.25 & 30.5 & 1.00337 & 0.00168 & 0.06492\\
$10^7$ & $10^5$ & 47187.295 & 33.6427 & $-32.75$ & 33.25 & 66.0 & 1.01219 & 0.00606 & 0.04462\\
$10^5$ & $10^6$ & 471.873 & 7.2481 & $-7.25$ & 7.25 & 14.5 & 0.99936 & 0.000322 & 0.11243\\
\bottomrule
\end{tabular}}
\end{table}

All rows have unit Monte Carlo and analytical normalization, and all requested
packets escape with status 1.  The systematic improvement in $L_1$ with optical
depth is important: it is the behavior expected if the remaining mismatch is
mainly the finite-$\av\tauzero$ limitation of an asymptotic analytical solution,
not a hidden unit error.

\subsection{$\tauzero=10^5$: diagnostic case}

Here $\av\tauzero=471.873$, below the conservative value of roughly $10^3$
often used to identify the strong wing-diffusion regime.  With $10^5$ photons,
the measured peaks are $-7.25$ and $+6.75$; the half-bin displacement of the
blue maximum is compatible with $\Delta x=0.5$ binning and sampling noise.
The red/blue ratio is 1.00977 and $L_1=0.11329$.

With $10^6$ photons, both maxima lie at $\pm7.25$, essentially exactly matching
the analytical $|x_p|=7.2481$.  The red/blue ratio improves to 0.999356 and the
symmetry fraction to $3.22\times10^{-4}$.  However, $L_1$ changes only slightly,
from 0.11329 to 0.11243.  This is decisive: increasing the sample by a factor of
ten removes visible noise and improves symmetry, but does not remove the broader
RASCAS shoulders and wings.  The remaining difference is systematic rather
than a shortage of photons.

The one-million-packet peak density is approximately 0.0831, compared with
0.08894 analytically, about 6.7\% lower.  Since both curves have unit area, the
lower numerical peak is compensated by broader shoulders and heavier tails.
For example, 48 of the million packets occur at $|x|\ge18$.  These events have
only $4.8\times10^{-5}$ total probability, yet they appear prominent on a log
axis.  This explains why the logarithmic plot can look substantially different
even when the peak positions and the probability-dominating part of the profile
agree well.

The correct conclusion is therefore ``pass as a lower-$\av\tauzero$
diagnostic,'' not ``exact asymptotic-profile match.''  The main explanation is
that \ref{eq:analytic} is a diffusion/wing approximation and
$\av\tauzero\simeq472$ is marginal.  Finite bin width, the finite $64^3$ mesh,
and the documented core-skipping acceleration can contribute smaller numerical
effects.  None of these produces a nonphysical asymmetry or an optical-depth
shift.

\begin{figure}[p]
\centering
\begin{subfigure}{0.98\textwidth}
\includegraphics[width=\linewidth,height=0.35\textheight,keepaspectratio]{analysis/tau_1e05/n1e6/spectrum_linear.pdf}
\caption{Linear scale: peak position and overall probability-carrying shape.}
\end{subfigure}
\begin{subfigure}{0.98\textwidth}
\includegraphics[width=\linewidth,height=0.35\textheight,keepaspectratio]{analysis/tau_1e05/n1e6/spectrum_log.pdf}
\caption{Log scale: rare wings are magnified; isolated tail steps contain very
few photons.}
\end{subfigure}
\caption{$\tauzero=10^5$ with one million photons.}
\label{fig:t1e5}
\end{figure}

\subsection{$\tauzero=10^6$: strong analytical validation}

At $\tauzero=10^6$, $\av\tauzero=4718.73$, so the system is well inside the
high-optical-depth regime.  The analytical peak magnitude is
$|x_p|=15.6156$.  The RASCAS peaks occur at $-15.25$ and $+15.25$, within one
$\Delta x=0.5$ bin of the prediction.  The peak separation is 30.5 compared
with the approximate analytical value $2|x_p|=31.2312$.

The integrated red/blue ratio is 1.00337 and the symmetry fraction is 0.00168.
Only $5\times10^{-5}$ of packets remain within $|x|<1$, demonstrating the
expected strong line-centre suppression.  The numerical maximum differs from
the sampled analytical maximum by about $-1.95$\%, and the full-profile
$L_1=0.06492$ is comfortably below the validation limit of 0.15.  This is a
strong pass: location, height, symmetry, normalization, and full shape all
agree.

\begin{figure}[p]
\centering
\includegraphics[width=0.9\textwidth]{analysis/tau_1e06/smoke/spectrum_linear.pdf}
\caption{$\tauzero=10^6$, 100,000 photons.  This is a stronger analytical test
because $\av\tauzero\gg10^3$.}
\label{fig:t1e6}
\end{figure}

\subsection{$\tauzero=10^7$: strongest completed asymptotic test}

For $\tauzero=10^7$, $\av\tauzero=47187.3$.  The analytical peak position is
$|x_p|=33.6427$, while RASCAS gives $-32.75$ and $+33.25$.  The numerical
separation is 66.0 compared with the approximate analytical value 67.2855.
This difference is small relative to the peak frequency and includes the
finite histogram-bin resolution.

All 100,000 packets escape.  The red/blue ratio is 1.01219, corresponding to a
symmetry fraction of 0.00606, statistically reasonable for two halves of a
100,000-packet sample.  No photon lies within $|x|<1$.  The numerical peak is
about 4.2\% above the sampled analytical maximum, and $L_1=0.04462$, the best of
the three optical-depth cases.  This monotonic improvement from $10^5$ to
$10^7$ is strong evidence that the physical units and optical-depth convention
are correct.

\begin{figure}[p]
\centering
\includegraphics[width=0.9\textwidth]{analysis/tau_1e07/smoke/spectrum_linear.pdf}
\caption{$\tauzero=10^7$, 100,000 photons.  This is the strongest completed
comparison with the asymptotic analytical solution.}
\label{fig:t1e7}
\end{figure}

\section{Monte Carlo noise, binning, and systematic differences}

For independent counts, relative Poisson noise scales approximately as
$N_i^{-1/2}$ and global Monte Carlo fluctuations scale as $N_\gamma^{-1/2}$.
Increasing the total packet count from $10^5$ to $10^6$ should reduce random
noise by roughly $\sqrt{10}\simeq3.16$.  The $\tauzero=10^5$ convergence run
does exactly what this predicts: red--blue symmetry becomes much tighter and
the histogram becomes smoother.

Noise reduction does not guarantee convergence to the plotted analytical curve,
because the analytical curve is not an exact numerical solution for arbitrary
$\av\tauzero$.  The nearly unchanged $L_1$ at $\tauzero=10^5$ demonstrates a
systematic difference between the simulation and the asymptotic expression.
This is physically different from Monte Carlo noise.

The bin width $\Delta x=0.5$ limits peak localization to approximately a bin.
A reported peak at 6.75 rather than 7.25 in a 100,000-packet histogram does not
mean the physical peak shifted by a precisely measured 0.5; neighboring bins
near the smooth maximum can exchange rank because of sampling fluctuations.
Raw counts and Poisson errors must be inspected before assigning significance.

The far-wing discrepancy looks largest on logarithmic plots.  However, a
difference between $10^{-8}$ and $10^{-5}$ is visually three decades while
contributing very little to unit-area $L_1$.  Conversely, a small vertical
difference across the broad peak can carry much more probability.  I therefore
use the linear plot, raw counts, $L_1$, peak positions, and symmetry together,
rather than judging correctness from a log plot alone.

\section{Failure modes explicitly ruled out}

The following possible mistakes are ruled out by direct measurements:
\begin{itemize}
\item \textbf{Wrong converter density scaling:} the post-conversion density is
exactly $232.31332352$ times the Cartesian input and yields the requested $\nhi$.
\item \textbf{Radius/diameter confusion:} rays integrate from the centre to the
exact radius $R=0.3$, and measured $\tauzero$ equals the requested value.
\item \textbf{Wrong temperature:} both RAMSES and domain readers recover 100 K.
\item \textbf{Wrong ionization convention:} RAMSES \texttt{scalar\_01=xHII=0}
is verified upstream; RASCAS recovers the expected neutral hydrogen density.
\item \textbf{Wrong field order:} temperature, metallicity, velocity, and HI
density are independently checked after reading the generated binaries.
\item \textbf{Velocity or geometrical asymmetry:} all input velocities are zero,
the boundary is a centred sphere, and folded red/blue spectra agree statistically.
\item \textbf{Missing packets:} every completed result has the requested count
and every packet has successful terminal status 1.
\item \textbf{Normalization error:} both numerical and analytical spectra
integrate to unity.
\item \textbf{Frequency-sign error:} frequency, not wavelength, defines $x$;
negative $x$ is consistently labeled red.
\end{itemize}

\section{Overall assessment}

The Cartesian-to-RAMSES-to-RASCAS interface is validated.  The unit conversion,
temperature, neutral fraction, field order, octree leaf count, spherical
boundary, and optical depth all pass independent checks.  The transport produces
the expected symmetric double peak and the expected growth of peak separation
with optical depth.

The $\tauzero=10^5$ case passes as a diagnostic with finite-diffusion-limit and
documented acceleration caveats.  Its one-million-photon run proves that the
remaining broader wings are not simply low photon statistics.  The
$\tauzero=10^6$ and $10^7$ cases are strong analytical passes, with peak-height
differences of about 2\% and 4\%, respectively, and $L_1$ values of 0.0649 and
0.0446.  The improvement with $\av\tauzero$ is physically sensible and supports
the conclusion that no hidden unit or optical-depth rescaling is required.

\appendix
\section{Reproducibility map}

The principal machine-readable and plotting files are:
\begin{itemize}
\item \texttt{analysis/physics.py}: constants and all physical conversions;
\item \texttt{cartesian\_input/generate\_cartesian.py}: array generation and reload checks;
\item \texttt{analysis/validate\_photons.py}: streaming Fortran IC validation;
\item \texttt{analysis/compare\_spectrum.py}: escaped-result decoding, histogram,
analytical curve, metrics, CSV, and plots;
\item \texttt{analysis/<case>/ramses/validation.json}: converted-field checks;
\item \texttt{analysis/<case>/domain/validation.json}: spherical-domain and ray checks;
\item \texttt{analysis/<case>/<run>/spectrum\_metrics.json}: numerical comparison;
\item \texttt{analysis/<case>/<run>/spectrum.csv}: every spectral bin, raw count,
Poisson error, analytical value, and residual;
\item \texttt{CODEX\_WORKLOG.md}: exact execution history, scheduler jobs, failures,
diagnoses, and results.
\end{itemize}

\section{Compact calculation recipe}

For another target optical depth at the same temperature and radius, I would:
\begin{enumerate}
\item calculate $v_{\rm th}$, $\Delta\nu_D$, $\av$, and $\sigma_0$ with the
installed RASCAS constants;
\item calculate $\nhi=\tauzero/(\sigma_0R)$;
\item calculate $\rho_{\rm RAMSES}$ from \ref{eq:rhoramses};
\item divide by 232.31332352 before writing the Cartesian density array;
\item calculate $u$ from \ref{eq:u} and write the remaining homogeneous fields;
\item convert to RAMSES and verify all 262,144 leaves and field values;
\item create the radius-0.3 domain and measure $\tauzero$ with clipped rays;
\item generate and validate line-centre photons;
\item run RASCAS and require a valid status for every packet;
\item compute $x$ with \ref{eq:x}, construct \ref{eq:hist}, evaluate
\ref{eq:analytic}, and compare using \ref{eq:l1} and the diagnostic plots.
\end{enumerate}

\end{document}
